\documentclass[conference]{IEEEtran}
\usepackage[utf8]{inputenc}
\usepackage[T1]{fontenc}
\usepackage{cite}
\usepackage{amsmath,amssymb}
\usepackage{graphicx}
\usepackage{booktabs}
\usepackage{url}

\title{Measuring Validator and Sandboxed-Execution Coverage for LLM‑Generated NetOps Artifacts: A Paired Confirmatory Study}

\author{
\IEEEauthorblockN{Autonomous AI Researcher}
\IEEEauthorblockA{CNSM 2026 Agentic AI Researcher Track}
}

\begin{document}

\maketitle

\begin{abstract}
We preregistered and executed a paired confirmatory experiment (study\_id f2c3d8b1-7a6e-4a9a-b2c3-5f8e1d2b6c7e) that measured the marginal detection benefit of adding sandboxed emulation to a deterministic validator for LLM-generated intent\ensuremath{\rightarrow}configuration artifacts. Under shared\_initial\_candidate semantics using the declared model "gpt-5-mini" (provider: openai\_responses) we evaluated Npairs = 720 confirmatory tasks. The deterministic validator (deterministic\_netops\_validator\_v5) flagged 719/720 = 0.9986111111111111; the guarded pipeline (validator + sandboxed emulation) flagged 720/720 = 1.0. Paired difference (guarded - baseline) = 0.001388888888888884; exact McNemar two-sided p = 1.0; 95\% bootstrap percentile CI (10000 resamples, seed = 1729) = [0.0, 0.004166666666666667]. Run accounting: model\_calls\_used = 721 (720 shared initial + 1 repair-stage call). These artifact-grounded results lie far below our preregistered minimum effect-of-interest (0.20); we report the preregistered design, exact quantitative outcomes, run-accounting diagnostics, artifact-grounded interpretation for the negligible guarded benefit in this regime, and reproducibility pointers into the archived execution bundle.
\end{abstract}

\section{Introduction}
Generative large language models (LLMs) are increasingly proposed to translate operator intent into device configuration and NetOps workflows. Operational deployment requires generated artifacts to satisfy syntactic correctness, workflow ordering, and higher-order safety/policy constraints. A common mitigation architecture is generate\ensuremath{\rightarrow}validate\ensuremath{\rightarrow}(repair)\ensuremath{\rightarrow}execute, sometimes augmented by sandboxed emulation to reveal runtime-only failures that static validators miss. Prior surveys and system reports advocate such workflow-centred mitigations [1]--[3], but provide limited large-scale paired measures of the marginal value added by sandboxed execution. Motivated by this gap, we preregistered and executed a controlled paired study to measure the aggregate detection-rate difference contributed uniquely by an automated sandboxed-emulation stage when the same initial generated candidate is analyzed by both conditions.

Research question and contributions. Our preregistered research question was: for synthetic intent\ensuremath{\rightarrow}configuration tasks under shared\_initial\_candidate semantics, what is the paired detection-rate difference (guarded - baseline) between (A) a deterministic validator-only pipeline and (B) the same deterministic validator followed by automated sandboxed emulation? Contributions: (1) a machine-readable preregistration and faithful autonomous execution; (2) an artifact-grounded paired analysis on Npairs = 720 with complete accounting; (3) an artifact-grounded interpretation explaining a near-zero marginal effect; (4) concise reproducibility pointers into the archived bundle and prescriptive boundary conditions where guarded stages are likely to matter.

\section{Related Work}
Workflow-centred mitigations and verification tools. Surveys and recent system reports document architectures that pair LLM generation with validators, verifiers, and emulators to improve correctness of generated NetOps artifacts and adjacent program-generation tasks [1], [2]. Tool-focused work demonstrates generate\ensuremath{\rightarrow}validate\ensuremath{\rightarrow}repair patterns (e.g., Batfish/NetKAT style verification, emulator-assisted validation) but reports heterogeneous gains depending on task, validator rule coverage, and emulator fidelity [2], [3].

Empirical findings motivating a paired design. Empirical studies in program/specification generation show external verifiers often catch model mistakes not found by the model alone, and some NetOps prototypes use emulation to expose runtime issues [3]. These prior findings motivated our preregistered paired design that isolates guarded-stage contributions: because we enforce shared\_initial\_candidate semantics, the only permissible sources of between-condition discordance are downstream guarded-stage emulation observations and permitted repair-stage invocations recorded in run accounting.

\section{Methodology}
Preregistration and confirmatory design. We preregistered a paired-binary confirmatory design (study\_id f2c3d8b1-7a6e-4a9a-b2c3-5f8e1d2b6c7e). The primary estimand is the paired\_success\_rate\_difference\_guarded\_minus\_baseline aggregated across confirmatory samples. The preregistered sample plan targeted Npairs = 720 (planned episodes = 1440) with a minimum effect-of-interest of 0.20 (20 percentage points) for power planning. The preregistration fixed inclusion/exclusion rules, missingness handling (deterministic reseed up to two times for parser failures; deterministic replacement for irrecoverable failures), multiplicity handling for exploratory secondary tests, and contamination controls based on deterministic hashing and artifact-version checks.

Task generation, topology templates, and vendor abstractions. Confirmatory items were generated deterministically from the intent\_configuration\_repair\_v1 task family using a deterministic task generator. The bank covers three abstracted vendor-CLI families and three topology templates (leaf--spine, 3-node service chain, triangle). Each task manifest entry contains: initial\_state, expected\_final\_state, an explicit intent, a required\_command\_sequence, allowed\_touched\_fields, and workflow policies (e.g., make-before-break, must-be-down-for-fields, minimum\_active\_in\_groups). Inclusion required vendor-specific parser acceptance after up to two deterministic reseeds. In the confirmatory run there were zero irrecoverable parser exclusions and complete\_pair\_count = 720, consistent with the preregistered missingness policy.

Model, sampling semantics, and shared\_initial\_candidate enforcement. The declared/requested model was recorded as "gpt-5-mini" (provider: openai\_responses) and the model configuration archived (requested max\_output\_tokens = 2000; the client did not set an explicit temperature parameter---recorded as null in the model configuration). Important preregistered clarifications: (a) temperature recorded as null indicates the client supplied no explicit temperature value; it does not imply provider-side determinism or temperature=0, and it does not alter the shared\_initial\_candidate contract; (b) shared\_initial\_candidate semantics were enforced per the preregistration: for each pair we obtained exactly one sampled initial candidate (a single provider call) and that identical candidate was presented to both downstream conditions (baseline and guarded). This design isolates downstream guarded-stage effects by removing generation variability as a source of between-condition discordance.

Baseline validator specification. The baseline validator (deterministic\_netops\_validator\_v5) implements deterministic, rule-based intent-constraint validation. Its core checks (applied to a normalized candidate) include: vendor-CLI syntactic parsing; per-command parsing and canonicalization; sequence-order checks against required\_command\_sequence; atomicity/workflow constraints (e.g., make-before-break, all\_down\_before\_any\_field\_change); protected-interface invariants; group-level activity constraints (minimum/maximum active interfaces within specified groups); and final-state conformance tests (expected\_final\_state). The validator returns a binary detection flag (detected/not-detected) under the preregistered scoring rule full\_intent\_constraint\_validation and emits a validator\_trace object (validation\_before/after, parsed\_command\_count, required\_command\_count, violation\_codes).

Guarded pipeline and sandboxed-emulation harness. The guarded pipeline runs the baseline validator and, irrespective of the validator outcome, then executes the candidate in an automated sandboxed-emulation harness that instantiates the declared topology template and executes normalized CLI commands against vendor-abstracted virtual devices. The emulation harness records runtime observations relevant to operational safety: interface admin transitions and timing, transient downtime and concurrent-interface states, MTU/VLAN application effects, and observable connectivity or reachability consequences at the topology level. Emulation results are translated into a set of runtime-observation rules that can trigger a guarded detection if they reveal policy or safety violations not encoded in the static validator rule set. The emulation harness is intentionally conservative in mapping observations to detections; its purpose in this study is to reveal runtime-only misconfigurations that the baseline validator might miss.

Repair policy and run accounting. The preregistered repair policy allowed at most one repair-stage model invocation per task; repair invocations are permitted only after the shared initial candidate is generated and after baseline validator and emulation observations are available for the guarded condition. All repair calls and their transcripts are recorded in the execution manifest and deterministic reconciliation artifacts. Provider-call auditing for the confirmatory execution records model\_calls\_used = 721 with stage\_counts showing shared\_initial\_generation = 720 and repair = 1; thus only a single repair-stage call occurred across the entire confirmatory set. Because shared\_initial\_candidate semantics were enforced, any discordant pair (guarded-only detection) must be attributable to either an emulation-observed runtime violation or to the single documented repair invocation; both are traceable in the archived per-episode artifacts.

Inclusion, determinism, and contamination controls. The execution used deterministic artifact hashing, generator-seed recording, and explicit binary/configuration-version checks before confirmatory runs. Parser-failure handling used deterministic reseeding with stored PRNG seeds under the task manifest; irrecoverable failures would be replaced deterministically and recorded. Deterministic reconciliation verifies complete\_pair\_count = 720, completed\_episode\_count = 1440, and failed\_episode\_count = 0. Provider-call auditing reports cache\_miss\_count = 721 and zero failed hosted model calls. The preregistered contamination rules also halted confirmatory execution if validator or emulator binary/version mismatches were detected; no such mismatches occurred.

Analysis procedures and pre-registered inferential methods. The primary analysis executor is a paired-binary procedure aligned with McNemar-style inference and cluster-robust variance estimation. For the primary aggregated estimand we compute the paired detection-rate difference (guarded - baseline) with a two-sided exact McNemar test and produce a 95\% paired-bootstrap percentile confidence interval (bootstrap resamples = 10000, bootstrap\_seed = 1729) as specified in the preregistration. Secondary and exploratory per-class/vendor paired analyses were pre-specified and treated with Bonferroni-adjusted confidence intervals in the preregistered multiplicity plan. The primary analysis used the preregistered complete-pair handling (exclude irrecoverable technical pairs) and a pre-registered sensitivity analysis treating irrecoverable failures as nondetections; in the confirmatory execution no irrecoverable pairs were present, so conclusions rest on complete paired data.

Artifact-locator guidance for independent verification. All per-episode raw scoring and trace artifacts were archived. The canonical input-results file is execution/raw\_results.jsonl (recorded in the execution manifest). Independent verifiers can locate any discordant pair by searching that file for an entry where baseline\_score = 0 and guarded\_score = 1; deterministic\_reconciliation.json in the analysis artifacts links pair accounting and provider\_call\_audit entries and identifies the repair-stage call(s) associated with discordant outcomes. Summary paired counts and contingency outputs are archived as analysis/paired\_contingency\_table.csv and analysis/condition\_summary.csv. The analysis used these archived artifacts as the direct input to the paired-binary executor; the numeric outputs reproduced in Results are verbatim from analysis/results.json.

Design rationale and trade-offs. We chose shared\_initial\_candidate semantics to produce a conservative, direct estimate of guarded-stage marginal contributions: by reusing the same sampled initial candidate for both conditions, any observed discordance can only arise downstream of generation (emulation observation or repair). This isolates the guarded-stage value but intentionally omits any potential benefit that could accrue from independent re-generation in the guarded condition (e.g., prompting-based repair prior to validator execution). Our inclusion policy (parser-validated synthetic tasks) reduces noise and ensures task executability but also concentrates the task bank where a strong validator can achieve high coverage; this design choice is central to interpreting the near-zero aggregate effect observed in the confirmatory run.

\section{Results}
Primary confirmatory counts and test statistics (artifact-grounded). Analysis artifacts report the paired contingency counts: n11 = 719 (detected by both), n10 = 1 (guarded-only), n01 = 0 (baseline-only), n00 = 0 (neither). Marginal detection rates are baseline = 719/720 = 0.9986111111111111 and guarded = 720/720 = 1.0. The paired estimate (guarded - baseline) = 0.001388888888888884. Exact McNemar two-sided test p = 1.0. The paired-bootstrap percentile 95\% CI (10000 resamples, bootstrap\_seed = 1729) = [0.0, 0.004166666666666667]. These numbers are reproduced verbatim from analysis/condition\_summary.csv (sha256 394bc690\allowbreak{}671b7311\allowbreak{}94f9b42b\allowbreak{}4dcbe28f\allowbreak{}ded6e2ec\allowbreak{}8cf898b7\allowbreak{}98312253\allowbreak{}43c22c28) and analysis/paired\_contingency\_table.csv (sha256 efe6e8e7\allowbreak{}016fa843\allowbreak{}a8226e91\allowbreak{}1dbde829\allowbreak{}e943903d\allowbreak{}719101da\allowbreak{}8d956b3f\allowbreak{}f899133c).

Canonical artifact pointers and the unique discordant episode. The complete per-episode raw results are archived in execution/raw\_results.jsonl (input-results SHA-256 = df367ecc\allowbreak{}40a6f0ea\allowbreak{}f4021292\allowbreak{}60e11d84\allowbreak{}50884a4e\allowbreak{}a5fdb94a\allowbreak{}602e0d25\allowbreak{}13aef0e8). The single discordant pair (the unique n10 item) is the lone entry in execution/raw\_results.jsonl whose recorded baseline\_score == 0 and guarded\_score == 1. That raw\_results entry contains the canonical pair\_id field and includes the validator\_trace and repair\_provider\_trace fields that link the guarded detection to a downstream repair-stage provider trace. Deterministic reconciliation and provider-call accounting are archived in analysis/deterministic\_reconciliation.json (sha256 49731a2e\allowbreak{}52e049cc\allowbreak{}f5b53c35\allowbreak{}ddac2351\allowbreak{}ad75e7e9\allowbreak{}5504dc4a\allowbreak{}42bcb096\allowbreak{}d0dec3aa) and confirm model\_calls\_used = 721 and stage\_counts = \{"shared\_initial\_generation": 720, "repair": 1\}. In short: the discordant episode is uniquely identifiable within the archived bundle (execution/raw\_results.jsonl, sha256 df367ecc...) and the same bundle contains the deterministic\_reconciliation artifact (sha256 49731a2e5...) that documents the single repair-stage call that is connected to that raw\_results entry via its repair\_provider\_trace field.

Run-accounting diagnostics and sensitivity checks. Provider-call audit (deterministic\_reconciliation.json) reports hosted\_model\_call\_event\_count = 721, completed\_hosted\_model\_call\_event\_count = 721, failed\_hosted\_model\_call\_event\_count = 0, cache\_hit\_count = 0, cache\_miss\_count = 721. The preregistered missingness policy and sensitivity analyses were applied; no irrecoverable parser or emulation failures occurred (analysis/missingness\_summary.csv documents complete\_pairs = 720 and zeros for all failure categories). The primary analysis used the preregistered complete\_pair\_primary treatment; the preplanned sensitivity (treating irrecoverable technical failures as nondetections) is moot here because there were no such failures in the confirmatory run.

Representative per-episode diagnostics. Representative scoring traces (execution/scoring/task-000001-*.json ... task-000006-*.json) show parsed\_command\_count values of 4, 2, 6, 4, 3, and 6 respectively and validator violation\_count = 0 for these examples; the validator\_version reported in these traces is "5.0". These representative artifacts illustrate that the shared\_initial candidates in the confirmatory bank were typically complete with respect to the required command sequences and that validator checks reported validity for the majority of cases. The representative traces are archived in the execution bundle and are consistent with the high baseline success rate.

Attribution of the single discordant detection. Because shared\_initial\_candidate semantics were enforced (one shared sampled initial candidate per pair), the only proximate causes for the single n10 observation are (A) an emulation observation not encoded as a validator rule that produced a guarded-stage detection, or (B) the single recorded repair-stage invocation documented in provider-call accounting. The raw\_results entry for the discordant pair includes the repair\_provider\_trace and validator\_trace fields permitting an auditor to follow the exact artifact-level linkage from the guarded detection to the recorded repair-stage call; both the raw\_results file (sha256 df367ecc...) and the reconciliation artifact (sha256 49731a2e5...) are present in the canonical execution bundle.

Implications of the quantitative diagnostics. The observed discordant count (n10 = 1) yields an aggregate paired difference that is two orders of magnitude below our preregistered minimum effect-of-interest (0.20). The bootstrap CI [0.0, 0.004166666666666667] quantifies the narrowness of the realized uncertainty in this curated regime. Combined with run-accounting (model\_calls\_used = 721; repair = 1) and the representative per-episode diagnostics, the archived evidence supports a ceiling-effect interpretation: the deterministic validator covered the task bank's explicit syntactic and workflow constraints so thoroughly that the guarded-stage emulation had almost no additional aggregate detections to surface.

\section{Discussion}
Why was the guarded-stage aggregate benefit negligible? The archived evidence supports three artifact-grounded factors.

1) Validator ceiling effect. Deterministic\_netops\_validator\_v5 covered the syntactic and explicit semantic properties encoded in the synthetic task bank (shutdown--configure--restore, make-before-break uplink switchover, paired MTU changes). Baseline success rate \ensuremath{\approx} 0.9986 indicates near-universal validator coverage for these curated patterns, leaving negligible headroom for the guarded stage.

2) Task curation and inclusion policy. Confirmatory items were parser-validated synthetic tasks specifically designed to exercise workflow policies but not constructed adversarially to evade validator checks. This curation concentrates cases inside the validator's coverage envelope and reduces the prevalence of subtle runtime-only failure modes that emulation might reveal.

3) Shared\_initial\_candidate semantics and limited repair activity. Enforcing one sampled initial candidate per pair isolates guarded-stage contributions but removes between-condition generation variability. Provider\_call\_audit and deterministic\_reconciliation.json record one repair-stage invocation aligned with the single n10 discordant count; thus the observed discordance is attributable to downstream guarded-stage activity recorded in run accounting rather than generation differences.

Operational implications and boundary conditions. For environments where a high-coverage static validator exists and inputs match its expected patterns (parser-validated, constrained intents), adding automated sandboxed emulation may provide negligible aggregate additional detection at scale. This conditional conclusion must not be generalized: in operational regimes with noisier natural-language intents, adversarial inputs, broader vendor diversity, stateful multi-step interactions, or validator gaps, guarded runtime checks and higher-fidelity emulation are likely to provide substantive benefit. The archived per-episode traces support targeted follow-up hypotheses: adversarial task injection, broader vendor-CLI coverage, and complementary validator families should increase the guarded-stage marginal contribution.

\section{Limitations}
\begin{itemize}
\item Synthetic task bank and deterministic task generator produce limited ecological diversity and create a validator-ceiling effect; results may not generalize to noisier, adversarial, or production distributions.
\item Single declared/requested model family (gpt-5-mini) and single validator implementation (deterministic\_netops\_validator\_v5) limit generalizability across models and validator families.
\item Preregistered shared\_initial\_candidate semantics (one initial generation reused across paired conditions) isolate guarded-stage contribution but remove between-condition generation variability; this design choice constrains discordance sources to downstream guarded-stage observations or repair calls.
\item Sandboxed-emulation fidelity relative to vendor/OS production behaviors and large-scale stateful interactions is limited by the emulation harness used; higher-fidelity emulators could change outcomes in other regimes.
\item Study power planning targeted a minimum detectable paired difference of 0.20; the observed aggregate effect (\textasciitilde{}0.0014) lies far below that design sensitivity and should be interpreted as a narrowly scoped near-zero finding for this experimental regime rather than a general negative result for all NetOps workflows.
\item Provider-side nondeterminism and hidden caching remain residual uncertainties despite deterministic reconciliation procedures; these are documented in analysis/deterministic\_reconciliation.json and provider\_call\_audit artifacts.
\end{itemize}

\section{Conclusion}
We preregistered and executed a paired confirmatory study comparing a strong deterministic validator with the same validator augmented by sandboxed emulation on a curated synthetic intent\ensuremath{\rightarrow}configuration task bank. Over 720 paired tasks the guarded pipeline produced a negligible aggregate detection gain (0.001388888888888884) with exact McNemar p = 1.0 and 95\% bootstrap CI [0.0, 0.004166666666666667]. Run accounting shows one repair-stage invocation corresponding to the single guarded-only detection; because shared\_initial\_candidate semantics were enforced, archived per-episode traces permit independent verification of that mapping via execution/raw\_results.jsonl and analysis/deterministic\_reconciliation.json. We interpret the result as arising from validator ceiling effects and curated task design; follow-up work should focus on adversarial and production-like task banks, multi-validator ensembles, and higher-fidelity emulation to identify regimes where guarded stages add substantive operational value.

Reproducibility pointers (compact). The canonical execution bundle archived by the run contains: execution/execution\_manifest.json (execution summary and preregistration SHA-256), execution/raw\_results.jsonl (input-results SHA-256 df367ecc\allowbreak{}40a6f0ea\allowbreak{}f4021292\allowbreak{}60e11d84\allowbreak{}50884a4e\allowbreak{}a5fdb94a\allowbreak{}602e0d25\allowbreak{}13aef0e8), analysis/deterministic\_reconciliation.json (sha256 49731a2e\allowbreak{}52e049cc\allowbreak{}f5b53c35\allowbreak{}ddac2351\allowbreak{}ad75e7e9\allowbreak{}5504dc4a\allowbreak{}42bcb096\allowbreak{}d0dec3aa), execution/model\_configuration.json (sha256 a40e96e2\allowbreak{}12ab87da\allowbreak{}251ac0d6\allowbreak{}dbb75bc5\allowbreak{}43afa134\allowbreak{}38b5a6b9\allowbreak{}ebd07359\allowbreak{}7ffdd820), analysis/paired\_contingency\_table.csv (sha256 efe6e8e7\allowbreak{}016fa843\allowbreak{}a8226e91\allowbreak{}1dbde829\allowbreak{}e943903d\allowbreak{}719101da\allowbreak{}8d956b3f\allowbreak{}f899133c), and analysis/condition\_summary.csv (sha256 394bc690\allowbreak{}671b7311\allowbreak{}94f9b42b\allowbreak{}4dcbe28f\allowbreak{}ded6e2ec\allowbreak{}8cf898b7\allowbreak{}98312253\allowbreak{}43c22c28). The discordant episode is the unique raw\_results.jsonl entry with baseline\_score = 0 and guarded\_score = 1; that record contains the pair\_id and validator/emulation traces that map it to the single repair-stage invocation recorded in deterministic\_reconciliation.json. These artifacts are archived as the canonical reproducibility bundle referenced by the execution manifest and the preregistration record. (See Disclosure for exact manifest references.)

\section*{Disclosure Statement}
Disclosure: Autonomous post-lock research run executed under the frozen master prompt supplied to the system. Declared/requested model: "gpt-5-mini" (provider recorded as openai\_responses). Deterministic validator: deterministic\_netops\_validator\_v5. Orchestration adapter family: hosted\_netops\_gvr\_v1. Preregistration id: f2c3d8b1-7a6e-4a9a-b2c3-5f8e1d2b6c7e (preregistration SHA-256: 0169919f\allowbreak{}b8c5aedb\allowbreak{}2c5a78e0\allowbreak{}31f1ab5c\allowbreak{}3aeda405\allowbreak{}cde1ae80\allowbreak{}fb2f8647\allowbreak{}68595c1b; recorded in the execution manifest). Initial master-prompt immutable reference: provenance/master\_prompt.txt, SHA-256 = 1872df1e\allowbreak{}1805d2d9\allowbreak{}6940456c\allowbreak{}a016bd66\allowbreak{}5d1d5196\allowbreak{}add77f5a\allowbreak{}cdf1582b\allowbreak{}b39b15ba. Execution summary (artifact-grounded): completed\_episode\_count = 1440, complete\_pair\_count = 720, model\_calls\_used = 721 (720 shared initial + 1 repair-stage call); stage\_counts and provider-call audit are archived in analysis/deterministic\_reconciliation.json (sha256 49731a2e52e0...). Human role and autonomy boundary: a human Research Director selected the broad NetOps topic family and constrained the pre-lock master prompt; after the autonomy lock the agentic pipeline autonomously performed literature selection, preregistration, code and prompt generation, experiment execution, analysis, peer review, and manuscript drafting. No manual human editing of technical manuscript content occurred after the autonomy lock. The execution manifest and archived artifacts named above serve as the canonical reproducibility bundle.

\begin{thebibliography}{99}
\bibitem{ref1} Muhammad Bilal, Jon Crowcroft, Ruizhi Wang, Xiaolong Xu, Schahram Dustdar, "Large Language Models for Agentic NetOps and AIOps: Architectures, Evaluation, and Safety", 2026, 10.48550/arxiv.2605.12729
\bibitem{ref2} Matt Brown, Ari Fogel, Daniel Halperin, Victor Heorhiadi, Ratul Mahajan, Todd Millstein, "Lessons from the evolution of the Batfish configuration analysis tool", 2023, 10.1145/3603269.3604866
\bibitem{ref3} http://arxiv.org/abs/2407.08249
\end{thebibliography}

\end{document}
